% page 181 of the facsimile (image p-180 of the scan)
% block 162.6 x 237.7 mm ; left margin 39.4 mm
\pgc{17.780mm}{0.000mm}{
\l{\cel{54}171}
\l{}
\l{\sou{folikulo}. I. (\sou{anat.)}\cel{2}Mikra sako membrana en la tegumento. -}
\l{\cel{3}II.\cel{1}\sou{(bo}t.Frukto kapsula, ek folio volvita longesale,ed}
\l{\cel{3}apertita per quaza suturo. - DEFIRS.}
\l{}
\l{\sou{\textquotedbl{}folium\textquotedbl{}}. (\sou{matem}.)Kurvo di qua parto esas pasable analoga a}
\l{\cel{3}la formo di arboro-folio.}
\l{}
\l{\sou{folkloro}. Civilizeso-formo materiala, strukturi sociala, tra-}
\l{\cel{3}dicionaji e literaturo-sorti parolala, metodologio (ca fako}
\l{\cel{3}inkluzante la tekniki yena : bibliografio, kartografio,}
\l{\cel{3}enrejistro sonala, inquesti, muzeografio). - DEFR.}
\l{}
\l{\sou{fomenta}r (\sou{trans.}) (\sou{medic.}) Remediar (parto di la korpo) per}
\l{\cel{3}aplikar e mantenar sur lu topiko varma. - DEFIS.}
\l{}
\l{\sou{fono}.\cel{2}(\sou{fiziko}).Un ek la serio de movi alternanta, tre rapida,}
\l{\cel{3}da korpo solida, liquido o gaso, quin la nervo audala}
\l{\cel{3}sentas per la orelo.}
\l{}
\l{\sou{fondar}. (\sou{trans.}) I. Apogar sur ulu, sur ulo (to quon onu entra-}
\l{\cel{3}prezas, esperas, kredas, e c.) . - II. Establisar (ulo) ye}
\l{\cel{3}maniero solida, ferma. - DEFIS.}
\l{}
\l{\sou{fonetiko}. La parto di linguistiko,qua traktas la son-uni (foni).}
\l{\cel{3}DEFIRS.}
\l{}
\l{\sou{fonetismo}. Reprezento exakta, per skribo, di la foni vocala.}
\l{}
\l{\sou{fongo}. (\sou{patol}.)Surkreskajo karnoza, sponjatra, quale fungo}
\l{\cel{3}(en ulcero, vunduri).}
\l{}
\l{\sou{fonografar}. (\sou{trans.}) Enrejistrar voco per aparato qua konsistas}
\l{\cel{3}ye disko vibranta, provizita ye stalo-stilo qua imprimas sur}
\l{\cel{3}cilindro o disko rotacanta la vibri recevita, tale ke esas}
\l{\cel{3}reproduktebla segun-vole, kande onu rotacigas la cilindro}
\l{\cel{3}o la disko, la vorti dicita avan la aparato. - DEFIRS.}
\l{}
\l{\sou{fonolito}. (\sou{mineral.).}\cel{2}Roko volkanala qua resonas sub martelo-}
\l{\cel{3}frapi. - DEF.}
\l{}
\l{\sou{fonometro}. (\sou{tekn.}) Aparato por mezurar la intenseso-grado di}
\l{\cel{3}la foni. - DEFIRS.}
\l{}
\l{\sou{fonometrio}. La arto mezurar la intenseso di la voco-foni.}
\l{}
\l{\sou{fonto}. I. Origino di aquofluo, di fonteno - aquo-veino qua pro-}
\l{\cel{3}duktas lu en la loko ube lu komencas ekirar de sulo. - II.}
\l{\cel{3}(metaf.) Origino dikozo, to de quo olca devenas, derivesas. -}
\l{\cel{3}EFIRS.}
\l{}
\l{fontanelo. (\sou{anat.})\cel{2}Che la infanteti, lakuno inter la osti kra-}
\l{\cel{3}niala, plenigita ye tisuo mola, flexebla, qua osteskas erste}
\l{\cel{3}kande la homo esas plu evoza. - DEFIS.}
}
